\documentclass[11pt]{article}

\usepackage{amsmath,amssymb}
\usepackage{xurl}
\usepackage[hidelinks]{hyperref}
\setlength{\emergencystretch}{2em}

% Shared commands for notes in docs/paper-gaps/.
% Notation follows the LDT paper (references/ldt-paper/).

\newcommand{\Lean}{\textsc{Lean}}
\newcommand{\leanid}[1]{\textsf{#1}}
\newcommand{\ghissue}[1]{\href{https://github.com/LionSR/MIPStarRE/issues/#1}{GitHub issue \##1}}
\newcommand{\ghpr}[1]{\href{https://github.com/LionSR/MIPStarRE/pull/#1}{GitHub PR \##1}}

% --- Paper-compatible notation ---
\newcommand{\F}{\mathbb{F}}
\newcommand{\N}{\mathbb{N}}
\newcommand{\C}{\mathbb{C}}
\newcommand{\tr}{\operatorname{tr}}
\newcommand{\ot}{\otimes}
\newcommand{\E}{\mathop{\bf E\/}}
\newcommand{\bx}{\mathbf{x}}
\newcommand{\by}{\mathbf{y}}
\newcommand{\bu}{\mathbf{u}}
\newcommand{\bv}{\mathbf{v}}

% Approximation relation (paper uses \approx_\eps and \simeq_\zeta)
\newcommand{\approxeps}{\approx_{\varepsilon}}
\newcommand{\simgeq}{\simeq_{\zeta}}

% Polynomial spaces (paper uses \polyfunc, \polysub, \polymeas)
\newcommand{\polyfunc}[3]{\operatorname{Poly}(#1,#2,#3)}
\newcommand{\polysub}[3]{\operatorname{PolySub}(#1,#2,#3)}
\newcommand{\polymeas}[3]{\operatorname{PolyMeas}(#1,#2,#3)}


\title{Issue 1093: The Total-Overlap Term in Projective Point-Consistency Transport}
\author{}
\date{2026-05-05}

\begin{document}
\maketitle

\section{The cited step}

This note concerns the projective point-consistency step in the
self-improvement theorem of \cite[Section~9]{Ji2020LowIndividualDegree}.%
\footnote{Tracked in \ghissue{1093} and \ghpr{1222}.  In the TeX source
consulted here, the projective point-consistency step is in
\path{references/ldt-paper/self_improvement.tex}, lines~719--726.  The
substitution proposition used there is \path{prop:triangle-sub} in
\path{references/ldt-paper/preliminaries.tex}.}
The helper stage produces a polynomial submeasurement $\widehat H$ which is
consistent with the point measurement $A$: on average over $u$,
\[
  A^u_a \ot I \simeq_{\widehat \zeta}
  I \ot \widehat H_{[h(u)=a]}.
\]
After orthonormalization, the paper obtains a projective submeasurement $H$ and
a postprocessed state-dependent distance estimate
\[
  I \ot \widehat H_{[h(u)=a]}
  \approx_{\widehat \zeta_{\mathrm{dataprocess}}}
  I \ot H_{[h(u)=a]}.
\]
The printed proof then applies \path{prop:triangle-sub} and concludes
\[
  A^u_a \ot I \simeq_{\widehat \zeta +
    \sqrt{\widehat \zeta_{\mathrm{dataprocess}}}}
  I \ot H_{[h(u)=a]}.
\]

\section{Where the total term enters}

In the measurement-valued right-register form of the triangle-substitution
argument, the consistency defect for a fixed question has the form
\[
  \max\left(0,\,
    \langle \psi, A_{\mathrm{tot}} \ot B_{\mathrm{tot}} \psi\rangle
    - \sum_a \langle \psi, A_a \ot B_a \psi\rangle\right).
\]
When $B$ is a measurement, $B_{\mathrm{tot}} = I$.  Thus replacing one
right-register measurement by another changes only the matching-mass term
\[
  \sum_a \langle \psi, A_a \ot B_a \psi\rangle.
\]
That term is exactly what the usual Cauchy--Schwarz estimate controls by the
square root of the state-dependent distance.

In the projective-output point-consistency step, however, the right-register
families are
\[
  u \longmapsto \widehat H_{[h(u)=a]},
  \qquad
  u \longmapsto H_{[h(u)=a]}.
\]
They are submeasurements, not measurements.  Their total operators are
$\widehat H_{\mathrm{tot}}$ and $H_{\mathrm{tot}}$, respectively, and these need
not both equal the identity.  Consequently the replacement also changes the
total-overlap term.  The additional scalar is
\[
  \mathbb E_u \left|
    \langle \psi, A^u_{\mathrm{tot}} \ot H_{\mathrm{tot}} \psi\rangle
    -
    \langle \psi, A^u_{\mathrm{tot}} \ot \widehat H_{\mathrm{tot}} \psi\rangle
  \right|.
\]
Since $A^u$ is a measurement, $A^u_{\mathrm{tot}} = I$, but this does not remove
the distinction between $H_{\mathrm{tot}}$ and $\widehat H_{\mathrm{tot}}$.

\section{Relation with the formal development}

The formal development therefore proves a submeasurement version of the
right-register substitution lemma with this total-overlap displacement stated
explicitly.%
\footnote{The preliminary statement is
\leanid{triangleSub\_right\_subMeas\_totalGap} in
\doclink{MIPStarRE/LDT/Preliminaries/Triangles.lean}.  The self-improvement
wrappers are \leanid{final\_fields\_point\_consistency\_totalGap\_natural} and
\leanid{final\_fields\_point\_consistency\_totalGap} in
\doclink{MIPStarRE/LDT/SelfImprovement/Theorems/Results/BoundednessTransport.lean}.}
In ordinary mathematical terms, the formal lemma says that if the helper
point-consistency error is at most $\delta$, the state-dependent distance
between the two right-register submeasurement families is at most $\varepsilon$,
and the averaged total-overlap displacement is at most $\eta$, then the
transported consistency error is at most
\[
  \delta + \sqrt{\varepsilon} + \eta.
\]
This is the same argument as the measurement-valued triangle-substitution
proposition, except that the term which vanishes for measurements is retained.

The formal statement does not yet prove an internal bound on $\eta$ from the
data-processing estimate.  A separate estimate may be possible, for example
with a dimension- or outcome-cardinality dependent loss obtained by applying
Cauchy--Schwarz to the sum of the outcome differences.  Such a bound is not the
literal measurement-valued step printed in the paper, and it has not yet been
formalized as part of the self-improvement theorem.

\section{Conclusion}

The printed proof applies a measurement-valued substitution principle at a
place where the formal objects are submeasurements.  The matching-mass part of
the argument is unchanged and is controlled by the square root of the
state-dependent distance, but a second term involving the right-register totals
appears.  This term is zero for measurements and is not zero in general for
submeasurements.

Thus the current formal development records a genuine proof-level gap in the
projective point-consistency paragraph.  The corrected formal statement includes
an additional total-overlap displacement parameter.  Closing the paper step as
stated requires either a proof that this displacement is absorbed by the
available data-processing error, possibly with a revised numerical constant, or
a different argument which avoids changing the total-overlap term.

\bibliographystyle{alpha}
\bibliography{references}

\end{document}
